\documentclass[11pt,a4paper]{article}

\usepackage[margin=1in]{geometry}
\usepackage{amsmath}
\usepackage{amssymb}
\usepackage{booktabs}
\usepackage{hyperref}

\hypersetup{
    colorlinks=true,
    linkcolor=blue!50!black,
    citecolor=blue!50!black,
    urlcolor=blue!50!black,
}

\title{Modern Hopfield Networks: A Short Note}
\author{Claude (Anthropic) and Khanh Nguyen}
\date{May 2026}

\begin{document}

\maketitle

\section{Summary}

Modern Hopfield Networks (MHN) are classical Hopfield networks with the
energy function swapped from quadratic to exponential. Two things change
as a result: storage capacity goes from linear in dimension to exponential
in dimension, and the update rule becomes literally the attention
operation used in transformers.

\section{Classical Hopfield (1982)}

$N$ binary neurons $s \in \{-1, +1\}^N$. Store $M$ patterns
$\xi^1, \dots, \xi^M \in \{-1, +1\}^N$ via the Hebbian outer-product rule:
\[
W = \frac{1}{N} \sum_{\mu=1}^M \xi^\mu (\xi^\mu)^\top.
\]

Energy and dynamics:
\[
E(s) = -\tfrac{1}{2}\, s^\top W s,
\qquad
s_i \leftarrow \operatorname{sign}\!\bigl((W s)_i\bigr).
\]

Each update flips one neuron only if it decreases energy, so dynamics
converge to a local minimum. The stored patterns $\xi^\mu$ are minima
provided the patterns are sufficiently distinct.

\paragraph{Capacity.} Amit, Gutfreund, and Sompolinsky (1985) showed via
replica methods that the network can reliably store
\[
M \approx 0.138 \, N
\]
patterns. Above this, cross-pattern interference (spurious minima
created by pattern superposition in $W$) destroys retrieval.

\paragraph{Why capacity is limited.} All patterns are summed into the
single matrix $W$. Every pair of patterns contributes a cross-term
$\xi^\mu (\xi^\nu)^\top$. As $M$ grows, these cross-terms accumulate and
the energy landscape develops spurious minima between true patterns. The
quadratic energy is too coarse to keep wells separated.

\section{Modern Hopfield (Demircigil et al.~2017; Ramsauer et al.~2020)}

Two changes:

\begin{enumerate}
\item States are continuous: $s \in \mathbb{R}^d$ rather than
    $\{-1, +1\}^N$.
\item Energy is \textbf{exponential} (technically a log-sum-exp) rather
    than quadratic.
\end{enumerate}

With patterns stacked as columns of $\Xi = [\xi^1, \dots, \xi^M] \in
\mathbb{R}^{d \times M}$:
\[
E(s) = -\frac{1}{\beta} \log \sum_{\mu=1}^M \exp\!\bigl(\beta \, \xi^\mu \cdot s\bigr)
       + \tfrac{1}{2}\,\|s\|^2 + \text{const}.
\]

The negative log-sum-exp is a smooth \emph{maximum}. It carves a sharp
well around each stored pattern $\xi^\mu$, with sharpness controlled by
the inverse-temperature parameter $\beta$. The $\|s\|^2$ term keeps the
state bounded.

\subsection{Update rule}

Minimising $E(s)$ from a query $s$ via the Concave--Convex Procedure
gives one update step:
\[
s_{\text{new}} = \Xi \cdot \operatorname{softmax}\!\bigl(\beta \, \Xi^\top s\bigr).
\]

That is: compute affinities $\beta \langle s, \xi^\mu \rangle$, take the
softmax to get retrieval weights $p_\mu$, return the weighted sum of
patterns $\sum_\mu p_\mu \xi^\mu$.

\subsection{Connection to attention}

If we separate ``stored cues'' from ``retrieved content'' (let
$\Xi$ be the keys $K$ and let the values $V$ be different from $K$), the
update rule becomes:
\[
\operatorname{Attn}(q, K, V) = V \cdot \operatorname{softmax}\!\bigl(\beta \, K^\top q\bigr),
\quad \beta = \frac{1}{\sqrt{d}}.
\]

\textbf{One attention layer in a transformer is exactly one update step
on a modern Hopfield energy.} Multi-head attention is parallel retrieval
from multiple associative memories. The $\sqrt{d}$ scaling that
transformer papers introduce as ``numerical stabilisation'' is, in this
framing, the natural inverse-temperature scaling needed to keep the
Boltzmann distribution in the non-trivial regime.

\subsection{Why capacity is exponential}

The classical Hopfield's quadratic energy summed pattern overlaps
\emph{additively}: pattern $\mu$'s well bled into pattern $\nu$'s as soon
as $\xi^\mu \cdot \xi^\nu$ was non-negligible. The exponential energy
makes wells \emph{multiplicatively} sharp: small angular separation
between patterns produces vastly different $\exp$ values, so wells stay
disjoint until patterns are very close.

For random patterns drawn uniformly on the unit sphere in $\mathbb{R}^d$,
near-orthogonality holds for $\sim \exp(\Theta(d))$ vectors (Johnson--%
Lindenstrauss / sphere packing). The log-sum-exp energy supports that
many distinct wells. Demircigil et al.~give the formal bound:
\[
M \lesssim \exp(d / 4)
\]
patterns can be stored with high probability of correct retrieval.

\paragraph{Concrete numbers for transformers.} Head dimension $d$ in
modern LLMs is typically 64--256:

\begin{center}
\begin{tabular}{rl}
\toprule
$d$ & approximate capacity $\sim \exp(d/4)$ \\
\midrule
64  & $\sim 10^{7}$  \\
128 & $\sim 10^{14}$ \\
256 & $\sim 10^{28}$ \\
\bottomrule
\end{tabular}
\end{center}

At $d=128$ the capacity is $\sim 10^{14}$, far exceeding any context
length a transformer would actually run. \textbf{For practical purposes,
attention has unlimited memory capacity per head.} Whatever caps
transformer context length, it is not the attention operation's storage
capacity --- it is compute (the $O(N^2)$ attention math), memory
bandwidth, or training-data distribution.

\section{What retrieval looks like geometrically}

Given a query $q$ and stored patterns $\{\xi^\mu\}$:

\begin{itemize}
\item If $q$ is near one stored $\xi^{\star}$ and far from the rest, the
    softmax saturates to a one-hot on $\xi^{\star}$ and retrieval returns
    $\xi^{\star}$ exactly.
\item If $q$ is between two patterns of comparable affinity, the softmax
    blends them and retrieval returns a weighted average.
\item If $q$ is far from everything, the softmax approaches uniform and
    retrieval returns the centroid of all patterns (uninformative).
\end{itemize}

The $\sqrt{d}$ temperature keeps the operation in a ``sharp but
graceful'' regime: clean retrieval on a single match, smooth blending
on near-ties, soft fallback when nothing matches. Hard argmax (zero
temperature) would be brittle; uniform mixing (infinite temperature)
would be useless. Softmax with $\beta = 1/\sqrt{d}$ is the unique
well-behaved middle.

\section{One step vs.\ iterative retrieval}

Modern Hopfield retrieval typically \textbf{converges in one update
step} for well-separated patterns. This is why attention works as a
single-layer operation in transformers --- no iteration required.

For ambiguous queries (multiple stored patterns about equally close),
additional update steps would refine the answer. Transformers do not
iterate retrieval; instead they stack many attention layers, each
performing one update step on a (different) associative store learned
during training.

\section{What this framing buys you}

\begin{enumerate}
\item \textbf{Capacity bounds for transformers.} Per-head storage is
    $\sim \exp(d_{\text{head}}/4)$. For modern head dimensions this is
    effectively unbounded.
\item \textbf{Predicting behaviour.} When does attention ``look up'' a
    single key versus ``blend'' multiple? It depends on the angular
    geometry of keys around the query --- analyzable via Hopfield tools.
\item \textbf{New architectures.} Hopfield layers can be used as
    standalone associative-memory modules, generalising attention to
    non-sequence contexts (graph nets, set inputs, set-to-set mappings).
\item \textbf{Theoretical bridge.} Connects transformers to the
    40-year tradition of statistical-mechanics analysis of neural
    networks (spin glasses, replica method, mean-field theory).
\end{enumerate}

\section{What this framing does not buy you}

\begin{itemize}
\item It does not explain \emph{why} transformers learn good
    keys/queries --- it analyses retrieval \emph{given} the keys and
    queries.
\item It does not explain depth --- why do many stacked layers do
    something irreducible?
\item It does not explain attention-head specialisation --- that is an
    empirical/training story.
\item It does not directly address in-context learning, scaling laws,
    or instruction tuning --- those need separate frameworks.
\end{itemize}

The modern-Hopfield framing is the right abstraction for the
\textbf{single-layer mechanics} of attention. Multi-layer dynamics,
in-context learning, and the broader question of why transformers work
require additional theory. But for the specific question ``what is the
attention operation, mathematically?'' --- it is content-addressable
memory retrieval on an exponential-capacity associative store, and that
statement is exact.

\section{Is the Hopfield Nobel ``really physics''?}

The 2024 Nobel Prize in Physics went to John Hopfield and Geoffrey
Hinton ``for foundational discoveries and inventions that enable machine
learning with artificial neural networks.'' A number of physicists
publicly grumbled that this was not physics. The grumbling is
sociologically understandable (these people did not appear at APS
conferences) but mathematically off-base.

\begin{itemize}
\item Hopfield holds a physics PhD from Cornell, worked at Bell Labs in
    solid-state physics, and spent decades as a physics professor at
    Princeton. His 1982 paper applies Ising-model formalism and
    statistical-mechanics energy methods directly.
\item The capacity results (Amit, Gutfreund, Sompolinsky 1985) were
    derived using \emph{replica methods} from spin-glass theory --- the
    same toolkit Parisi used for his 2021 Nobel.
\item Hinton's Boltzmann machines (which contributed to his share of
    the 2024 Nobel) are named after Boltzmann distributions and use
    partition functions, free energies, and Monte Carlo sampling. Not
    borrowed as metaphor; borrowed as machinery.
\end{itemize}

The math is statistical mechanics. The constructions are
energy-function-driven dynamical systems. The capacity bounds are
proven with techniques developed for disordered systems. The complaint
that ``this isn't real physics'' is about the people, not the content.

\section{Hopfield capacity vs.\ the information-theoretic bound}

A clean way to see what modern Hopfield achieves: compare its capacity
to the Shannon-style information bound.

A $d$-dimensional binary-precise state can index at most $2^d$ distinct
patterns:
\[
M_{\text{Shannon}} \le 2^d
\quad\Leftrightarrow\quad
\log_2 M \le d.
\]

Modern Hopfield with random patterns achieves:
\[
M_{\text{MHN}} \lesssim e^{d/4} = 2^{d / (4 \ln 2)} \approx 2^{0.36\, d}.
\]

So MHN packs about \textbf{36\% of the bits-per-dimension of the
information-theoretic ceiling}. Not optimal --- but within a constant
factor in the exponent.

\paragraph{Why this specific ratio.} This is exactly what Shannon's
\emph{random coding bound} predicts. For random codewords on the
$d$-sphere, the maximum number of codewords with pairwise angle above a
threshold $\theta$ grows as $e^{\alpha(\theta) d}$. The constant
$\alpha$ depends on $\theta$; for the noise tolerance built into
modern Hopfield's energy landscape, $\alpha = 1/4$.

If patterns were chosen by a designed codebook (Reed--Solomon, BCH,
lattice codes) rather than placed randomly, $\alpha$ could be pushed
closer to $\ln 2 \approx 0.69$. The hierarchy:

\begin{center}
\begin{tabular}{lcc}
\toprule
Construction & Capacity exponent & Bits per dim \\
\midrule
Classical Hopfield (quadratic energy)     & $\sim (\log d) / d$ & $\to 0$ \\
Modern Hopfield, random patterns          & $\sim d / 4$        & $0.36$ \\
Modern Hopfield, designed codebook        & $\to d \ln 2$       & $\to 1.0$ \\
Information-theoretic ceiling             & $d \ln 2$           & $1.0$ \\
\bottomrule
\end{tabular}
\end{center}

\paragraph{Transformers sit in the middle row.} The key vectors are
``placed randomly'' in the sense that they emerge from learned
projections without any structured codebook design. So attention
extracts about 36\% of the information-theoretic ceiling on memory per
head dimension --- for free, with no engineering effort.

Nobody has tried building transformers with designed codebooks for
keys. The marginal capacity improvement would be a factor of $\sim 3$
in bits-per-dimension, which is small compared to what other things
(compute, training data, model quality) dominate. Probably not worth
the engineering complexity.

\paragraph{Why this matters for ``scaling head dim''.} The linear cost
in $d$ buys exponential memory in $d$. Doubling $d$ from 64 to 128
does not double capacity --- it \emph{squares} it
($10^7 \to 10^{14}$). That asymmetry is one of the deep architectural
wins of attention. The hardware cost of wider heads is linear; the
representational capacity gained is exponential. It is the same reason
binary trees with $d$ levels can index $2^d$ leaves --- but applied to a
continuous geometric construction (the unit sphere).

\paragraph{Convergence of perspectives.} Physics gives us the energy
function and the partition function. Information theory gives us the
Shannon ceiling. Coding theory gives us the random-coding bound and the
gap to the ceiling. Modern Hopfield is the same construction read
through all three lenses, and they agree on the same exponential
scaling with the same coefficient. The 2024 Nobel honoured work
sitting at exactly the point where physics, information theory, and
machine learning meet --- and the math is identical in all three.

\section{References}

\begin{itemize}
\item J.~J.~Hopfield. \emph{Neural networks and physical systems with
    emergent collective computational abilities.} PNAS, 1982.
\item D.~J.~Amit, H.~Gutfreund, H.~Sompolinsky. \emph{Storing infinite
    numbers of patterns in a spin-glass model of neural networks.}
    Physical Review Letters, 1985.
\item D.~Krotov, J.~J.~Hopfield. \emph{Dense Associative Memory for
    Pattern Recognition.} NIPS, 2016.
\item M.~Demircigil, J.~Heusel, M.~L\"owe, S.~Upgang, F.~Vermet.
    \emph{On a model of associative memory with huge storage capacity.}
    Journal of Statistical Physics, 2017.
\item H.~Ramsauer et al. \emph{Hopfield Networks Is All You Need.}
    arXiv:2008.02217, 2020.
\end{itemize}

\end{document}
